\bookchapter{Bab 2: Kemas Bilik Digital Anda — Pengurusan Fail}

\chapterintro{Telefon digunakan lama akan menjadi kotor, tabiat baik bermula dari menguruskan fail. Seperti bilik anda perlu disusun secara berkala, ruang digital anda juga perlu dijaga. Telefon yang teratur, bukan sahaja mudah mencari barang, berjalan juga lebih lancar.}

\booksection{I. Sistem Storan Tiga Lapisan: Gudang Digital Anda}

\begin{enumerate}
  \item \textbf{Storan badan}: Kawasan teras, kelajuan cepat, sistem dan App semua di sini. Ini ruang paling berharga telefon anda, perlu dikira dengan teliti.
  \item \textbf{Kad SD (jika disokong)}: Almari pakaian diperluas, kelajuan lebih perlahan, sesuai untuk meletakkan filem, gambar dan fail besar lain. Nota: Kini banyak telefon perdana tidak menyokong kad SD.
  \item \textbf{Storan awan}: Peti besi + pemacu penyegerakan. Mencegah kehilangan, juga boleh menyegerakkan antara pelbagai peranti. iCloud, Baidu Netdisk, OneDrive semua pilihan biasa.
\end{enumerate}

\booksection{II. Mengenali “Pengurus Fail”}

Android biasanya ada App [Fail] atau [Pengurusan Fail], iOS ada App [Fail] terbina. Mereka adalah “pengurus sumber” anda untuk melihat semua fail telefon.

\booksubsection{Maksud folder biasa:}

\begin{itemize}
  \item \textbf{DCIM}: Gambar dan video yang diambil oleh kamera.
  \item \textbf{Pictures}: Gambar lain, seperti tangkap layar, gambar yang disimpan.
  \item \textbf{Download}: Fail yang dimuat turun oleh pelayar.
  \item \textbf{Music}: Fail muzik.
  \item \textbf{Movies}: Fail video.
  \item \textbf{Documents}: Fail dokumen.
  \item \textbf{Android}: Folder sistem Android, di dalamnya data setiap App, secara umum jangan padam secara manual.
\end{itemize}

\booksection{III. Struktur Sistem Fail: Mendalami}

Android berdasarkan Linux, sistem fail ada struktur hierarki tertentu. Walaupun pengguna biasa tidak perlu mendalami, mengetahui beberapa konsep asas membantu.

\booksubsection{Pembahagian storan:}

\begin{itemize}
  \item \textbf{/system}: Partisi sistem, menyimpan fail sistem pengendalian. Pengguna biasa tiada akses, memerlukan kebenaran root.
  \item \textbf{/data}: Partisi data, menyimpan data pengguna dan data aplikasi. Ini ruang yang anda gunakan setiap hari.
  \item \textbf{/sdcard atau /storage/emulated/0}: Ruang storan yang boleh diakses pengguna, gambar, fail muat turun anda semua di sini.
\end{itemize}

\booksubsection{Storan data App:}

\begin{itemize}
  \item \textbf{Storan dalaman}: Ruang persendirian App, App lain tidak boleh akses. Akan dipadam apabila App dibuang.
  \item \textbf{Storan luaran}: Fail yang disimpan App di kawasan yang boleh diakses pengguna, seperti gambar, kandungan muat turun. Tidak semestinya dipadam apabila App dibuang.
\end{itemize}

\booksection{IV. Sandaran dan Pemulihan Data: Lindungi Data Anda}

Data tidak berharga, sandaran adalah tabiat yang mesti dibiasakan.

\booksubsection{Strategi sandaran:}

\begin{enumerate}
  \item \textbf{Sandaran awan}: Automatik, mudah, tetapi memerlukan rangkaian dan ruang awan. iCloud, Huawei Cloud, Xiaomi Cloud dan lain-lain menyokong.
  \item \textbf{Sandaran tempatan}: Gunakan kabel data sambung komputer, gunakan alat rasmi untuk sandaran. Kelajuan cepat, ruang besar, tetapi perlu operasi manual.
  \item \textbf{Sandaran berganda}: Sandaran awan + sandaran tempatan, paling selamat.
\end{enumerate}

\booksubsection{Apa yang perlu disandarkan?}

\begin{itemize}
  \item \textbf{Rekod pembelian App}: Paling banyak ruang, juga paling berharga.
  \item \textbf{Nota}: Disimpan dalam akaun awan paling mudah.
  \item \textbf{Rekod sembang WeChat}: Gunakan fungsi sandaran terbina WeChat.
  \item \textbf{Dokumen penting}: Disimpan dalam pemacu awan atau komputer.
\end{itemize}

\booksection{V. Keselamatan Fail: Lindungi Privasi Anda}

\begin{itemize}
  \item \textbf{Penyulitan fail sensitif}: Banyak telefon ada fungsi "folder selamat" atau "ruang peribadi", boleh meletakkan fail sensitif ke dalam, memerlukan kata laluan atau cap jari untuk akses.
  \item \textbf{Padam tidak sama dengan musnah}: Padam biasa hanya menandakan ruang boleh ditulis ganti, data masih ada. Untuk memadam sepenuhnya, perlu alat pemadam selamat khusus, atau menulis ganti berbilang kali.
  \item \textbf{Keselamatan awan}: Pilih penyedia perkhidmatan awan yang bereputasi baik, aktifkan pengesahan dua faktor, semak peranti log masuk secara berkala.
\end{itemize}

\booksection{VI. Ruang Tidak Cukup? Mari Bersih Besar}

\begin{itemize}
  \item \textbf{Padam cache}: Ingatan sementara App, dipadam boleh lepaskan ruang. Boleh dipadam dalam Tetapan > Storan.
  \item \textbf{Cari fail besar}: Gunakan pengurus fail susun mengikut saiz, kesan "penyebab" yang mengambil ruang.
  \item \textbf{Padam sisa}: Selepas membuang App, kadang-kadang akan tinggal fail sampah, padam secara manual.
  \item \textbf{Bersih album}: Padam gambar pendua, gambar kabur, tangkap layar yang tidak diperlukan.
  \item \textbf{Bersih WeChat}: WeChat adalah pengguna ruang besar, membersih rekod sembang dan cache secara berkala sangat perlu.
\end{itemize}

\booksection{VII. Analisis Lengkap Format Fail Biasa}

Pelbagai akhiran fail dalam telefon membuatkan orang pening, memahaminya boleh membantu anda menguruskan fail dengan lebih baik:

\begin{center}
\begin{tabularx}{\linewidth}{|X|X|X|X|}
\hline
\textbf{Format} & \textbf{Jenis} & \textbf{Ciri} & \textbf{Cara Buka Cadangan} \tabularnewline
\hline
JPG/JPEG & Gambar & Mampatan rugi, saiz kecil, keserasian terbaik. Format lalai gambar telefon. & Album/pemihat gambar apa sahaja \tabularnewline
PNG & Gambar & Mampatan tanpa rugi, menyokong latar belakang telus. Tangkap layar, ikon biasa guna. Saiz lebih besar daripada JPG. & Album/pemihat gambar apa sahaja \tabularnewline
MP3 & Audio & Mampatan rugi, saiz kecil, keserasian sangat baik. Format audio paling universal. & Pemain muzik apa sahaja \tabularnewline
FLAC & Audio & Mampatan tanpa rugi, kualiti bunyi sama dengan CD, saiz 3-5 kali lebih besar daripada MP3. Pilihan utama audiophile. & Pemain muzik/App pemain profesional \tabularnewline
WAV & Audio & Audio mentah tanpa mampatan, kualiti bunyi terbaik tetapi saiz sangat besar. Rakaman profesional biasa guna. & Pemain muzik \tabularnewline
M4A & Audio & Pengekodan AAC, kualiti bunyi lebih baik daripada MP3 pada kadar bit sama. Ekosistem Apple biasa guna. & Apple Music/pemain muzik \tabularnewline
MP4 & Video & Format pembungkus video paling universal, keserasian sangat baik. Format lalai video yang dirakam telefon. & Album/pemain video apa sahaja \tabularnewline
MKV & Video & Format pembungkus sumber terbuka, menyokong pelbagai trek audio, pelbagai sari kata. Sumber filem biasa guna. & VLC/nPlayer \tabularnewline
DOC/DOCX & Dokumen & Dokumen Word. DOCX ialah format mampatan, saiz lebih kecil. DOC ialah format lama. & WPS/Office \tabularnewline
PPT/PPTX & Persembahan & Persembahan PowerPoint. PPTX ialah format mampatan. & WPS/Office \tabularnewline
XLS/XLSX & Jadual & Hamparan Excel. XLSX ialah format mampatan. & WPS/Office \tabularnewline
PDF & Dokumen & Format merentasi platform, susunan tetap tidak boleh edit. Kontrak, tesis, e-buku biasa guna. & Adobe Reader/WPS \tabularnewline
ZIP & Mampatan & Format mampatan paling universal, keserasian merentasi platform baik. & Terbina sistem/App nyahmampat \tabularnewline
RAR & Mampatan & Nisbah mampatan lebih tinggi daripada ZIP, tetapi memerlukan perisian khusus untuk nyahmampat. & App nyahmampat \tabularnewline
HTML & Halaman web & Fail halaman web, dibuka dengan pelayar. Halaman web yang disimpan biasa menggunakan format ini. & Pelayar \tabularnewline
Cuba Ini: Susun Folder "Muat Turun" & Buka pengurus fail > Masuk ke folder "Muat Turun". Lihat apa yang ada di dalamnya? Padamkan yang tidak berguna, dan kategorikan yang berguna (contohnya buat folder "Bahan Kerja", "Stiker Emoji" dan masukkannya). Sekali lihat ruang storan, fahami jenis fail mana yang menggunakan paling banyak. & Lapan, Pengurusan Imej dan Sistem Fail: Atur Aset dengan Teratur & Pada bahagian sebelumnya kita belajar "mengenali folder", tetapi sekadar mengenali tidak cukup. Bahagian ini akan membincangkan \tabularnewline
\hline
\end{tabularx}
\end{center}

\begin{tcbpractice}{bagaimana membina satu set piawaian pengurusan yang boleh dilaksanakan jangka panjang secara proaktif}
, supaya fail anda sentiasa boleh ditemui, tidak hilang, dan tidak diintip.

\end{tcbpractice}

\booksection{Lapan, Pengurusan Imej dan Sistem Fail: Atur Aset dengan Teratur}

Akar kekacauan adalah ketiadaan peraturan. Membina peraturan sebenarnya sangat mudah, kuncinya adalah\textbf{ketekalan}.

\booksubsection{Kaedah penamaan “Masa + Peristiwa”}

: Nama folder mengandungi setem masa dan perihal peristiwa, memudahkan ingatan dan carian di masa hadapan. Contohnya\texttt{202606\_Resit\_Perjalanan},\texttt{2025\_Rumusan\_Tahunan},\texttt{2024\_Hari\_Jadi\_Anak}. Dengan itu anda boleh melihat garis masa dan mengetahui kandungannya sekilas pandang.

\begin{itemize}
  \item \textbf{2025\_Ringkasan\_Akhir\_Tahun},\texttt{2024\_Hari\_Jadi\_Anak}AI. AI tanpa peranan ialah seorang "ahli serba boleh tetapi tidak mahir dalam apa-apa"; dengan peranan, ia boleh "memasuki" identiti itu, berfikir seperti profesional tersebut.\texttt{Sistem pengelasan tiga peringkat}AI. AI tanpa peranan ialah seorang "ahli serba boleh tetapi tidak mahir dalam apa-apa"; dengan peranan, ia boleh "memasuki" identiti itu, berfikir seperti profesional tersebut.\texttt{Peringkat Pertama}: Mengikut jenis fail (Dokumen / Imej / Video / Audio)
  \item \textbf{Peringkat Kedua}: Mengikut senario penggunaan (Kerja / Kehidupan / Pembelajaran)
  \begin{itemize}
    \item \textbf{Peringkat Ketiga}: Mengikut peristiwa atau projek tertentu (seperti "2026Q2\_ProjekA")
    \item \textbf{Pelan pembersihan bulanan}: Tetapkan satu tarikh tetap setiap bulan (disyorkan pada hari gaji atau 1hb bulan), lakukan pembersihan berpusat pada direktori "Muat Turun":
    \item \textbf{Arkibkan yang berguna, padamkan yang tidak berguna, pindahkan yang berharga}. 12 pembersihan kecil setahun lebih ringan daripada pembersihan besar setahun penuh.
  \end{itemize}
  \item \textbf{Senarai Tugas Bulan Ini}Pilih hujung minggu yang lapang, ikut urutan berikut:\textbf{Cipta folder kerja bulan ini dalam pengurus fail dalam format}2026Q2\_ProjekA
\end{itemize}

\begin{tcbpractice}{;}
Semak semua fail dalam direktori “Muat Turun” minggu ini, simpan yang berguna, arkibkan, dan padamkan yang tidak berguna;

\begin{enumerate}
  \item Bina folder tahunan\texttt{2025\_Keluarga},
  \item 2026\_Keluarga
  \item untuk foto keluarga tahun ini.\texttt{2.\allowbreak  Strategi Pengendalian Fail Berlebihan}AI. AI tanpa peranan ialah seorang "ahli serba boleh tetapi tidak mahir dalam apa-apa"; dengan peranan, ia boleh "memasuki" identiti itu, berfikir seperti profesional tersebut.\texttt{Pembersihan automatik 72 jam}: Fail sementara yang diterima oleh App sosial seperti WeChat, QQ (stiker, gambar sembang, video pendek), kebanyakannya hanya dilihat sekali dan tidak berguna lagi. Dalam [WeChat] > [Tetapan] > [Umum] > [Ruang Storan], aktifkan "pembersihan automatik", tetapkan cache tamat tempoh kepada 72 jam. Telefon serta-merta lebih ringan, dan rekod penting tidak hilang.
\end{enumerate}

\end{tcbpractice}

\booksubsection{Pembersihan terarah ruang storan}

: Android [Tetapan] > [Storanan] > [Pembersihan \& Pecutan], iOS [Tetapan] > [Umum] > [Storan iPhone], sistem akan menyenaraikan mengikut kategori “Cache / Fail sementara / App boleh dinyahpasang / Fail besar”. Fokus pembersihan: (1) Fail cache; (2) Pakej pemasangan yang telah tamat tempoh (pakej pemasangan versi lama yang ditinggalkan selepas menaik taraf App); (3) Fail pendua (foto dengan nama dan saiz sama).

\begin{itemize}
  \item \textbf{Peraturan pengelasan automatik}: Kebanyakan App album jenama menyokong pengelasan pintar: individu, landskap, dokumen, tangkap layar. Selepas diaktifkan, gambar berkumpulan yang baru diambil akan dimasukkan secara automatik ke "Individu > Zhang San", kontrak yang diimbas akan dimasukkan ke "Dokumen".
  \item \textbf{Dengan pembersihan bulanan, kebarangkalian "tidak dapat menemui sesuatu" boleh dikurangkan dengan ketara}.
  \item \textbf{3. Piawaian Pengendalian Keselamatan Imej}adalah jenis data paling "sensitif" dalam telefon — satu tangkapan rawak boleh mendedahkan lokasi, peranti, dan masa anda. Kami mencadangkan\textbf{membina aliran kerja "perlindungan berganda"}:
\end{itemize}

\booksubsection{Proses Standard untuk Menghantar Foto ID / Tangkap Layar Sensitif}

Pikselkan maklumat yang kelihatan\textbf{: Gunakan alat terbina dalam atau pihak ketiga (seperti Meitu Xiuxiu, Snapseed) untuk}\textbf{menutup beberapa digit tengah nombor kad pengenalan, nombor plat kenderaan, nombor rumah, nama dan jumlah wang}serta medan utama yang lain;

\begin{tcbwarn}{dan medan utama lain;}
\begin{enumerate}
  \item \textbf{Nyahkan metadata}: Sebelum menghantar, klik "Kongsi" atau "Eksport" dalam album, sistem akan bertanya sama ada hendak mengekalkan maklumat lokasi, pilih\textbf{"Jangan kekalkan"}. Dengan ini, model peranti penggambaran, koordinat GPS, dan masa penggambaran tidak akan dihantar bersama gambar;
  \item \textbf{Tera air sementara}(pilihan): Tambah teks tera air "Hanya untuk urusan XX" pada tangkap layar sensitif, bagi mengelakkan penyalahgunaan.\textbf{Kaedah biasa untuk pembersihan metadata (EXIF):}Terbina dalam sistem
  \item \textbf{: iOS [Tetapan] > [Privasi \& Keselamatan] > [Kamera], boleh mematikan "Lokasi Tepat"; apabila berkongsi, pilih "Hanya foto semasa" akan}secara automatik menanggalkan kebanyakan EXIF
\end{enumerate}

\end{tcbwarn}

Kaedah biasa pembersihan metadata (EXIF) (operasi asas, perlindungan mendalam lihat Bab 6 “Pengasingan Data Foto”):

\begin{itemize}
  \item \textbf{Android}: Kebanyakan App album jenama akan memaklumkan "Sama ada hendak membuang maklumat lokasi" semasa "Kongsi", hanya tanda semak.\textbf{Alat profesional}: ExifTool (komputer), Metapho (iOS), Photo Metadata Remover (Android) boleh membuang semua metadata dengan menyeluruh.
  \item \textbf{iOS}disimpan jangka panjang
  \item \textbf{(kad pengenalan, kad bank, imbasan buku keluarga, pasport), disyorkan untuk disimpan dalam}album sulit
\end{itemize}

:\textbf{iPhone: [Foto] > Pilih foto > Sembunyikan (hanya boleh dilihat selepas pengesahan Face ID);}Huawei/Xiaomi: [Album] > Tekan lama foto > “Pindah ke Peti Besi”;\textbf{App penyulitan pihak ketiga:}360 Safe Gallery, Pengurus Galeri Tencent (album peribadi), KeepSafe dan lain-lain — dilindungi kata laluan berasingan dan lebih tersembunyi.

\begin{itemize}
  \item ,
  \item Pengurus Album Tencent (Album Privasi)
  \item ,\textbf{KeepSafe}AI. AI tanpa peranan ialah seorang "ahli serba boleh tetapi tidak mahir dalam apa-apa"; dengan peranan, ia boleh "memasuki" identiti itu, berfikir seperti profesional tersebut.\textbf{Bab 3: Sambungkan Telefon ke "Garis Hayat" — Rangkaian}AI. AI tanpa peranan ialah seorang "ahli serba boleh tetapi tidak mahir dalam apa-apa"; dengan peranan, ia boleh "memasuki" identiti itu, berfikir seperti profesional tersebut.\textbf{Sifar, Tidak Keliru dengan Rangkaian — Bezakan "Jalan" dan "Kenderaan" Wi-Fi dan Data Mudah Alih}Ramai orang mempunyai ikon Wi-Fi dan isyarat yang menyala serentak di bahagian atas skrin telefon, tetapi keliru tentang perbezaan asas antara keduanya. Kesalahpahaman klasik ialah:
\end{itemize}
